\section{Overview of lattice study.}

We use an ordinary single plaquette Wilson action defined on a
hypercubic lattice. We exploit twisted reduction \cite{twist8}
which allows a decent interval of lattice spacings to be realized
on a small lattice of size $V=7^4$.
For twisting to work
best, we need $N$ to be given by $4$ times a prime number.
We work mostly at $N=44$ with
checks for finite $N$ effects at $N=28$.
Our simulation method
employs a combination of heat-bath and over-relaxation updates and
equilibration is achieved in reasonable lengths of computer time.
In most of our analysis we can ignore all statistical errors.
As with any lattice work with few
precedents to rely on, a substantial part of the effort is to
identify useful ranges of the many parameters that enter numerical
work, so that the relevant physics questions can be addressed.
Throughout, we use $b$ for the lattice gauge coupling which is the
inverse bare 't Hooft coupling; relatively to the conventional
lattice coupling $\beta$, $b$ already has the right power of $N$
extracted:
\be
b=\frac{\beta}{2N^2}=\frac{1}{g_{\rm YM}^2N}\ee
For clarity purposes, we ignore a technicality in our
presentation:
For a twisted simulation $b$ is taken as the
negative of its value in the untwisted case; we shall quote the
conventional, positive value of $b$.

\subsection{Smeared lattice Wilson loops.}

We consider only rectangular loops restricted to one plane of our
hypercubic lattice and most of our work is done with square loops.
However, a continuum approach would
be concerned with the class of
observables that contain smooth loops of arbitrary shape.
Specifically, explicit calculations in the continuum would best be applied to
circular loops contained in one flat plane in $R^4$, as they have
maximal symmetry.


In order to eliminate the perimeter and corner
divergences \cite{div9} we construct the Wilson loop operators
associated with our curves not from the $SU(N)$ link matrices
$U_\mu (x)$ generated by the probability distribution determined
by the single plaquette Wilson action, but, rather, from smeared
versions of those. We employ APE smearing \cite{ape10}, defined
recursively from $n=0$, where the smeared matrix is equal to
$U_\mu(x)$. Let $\Sigma_{U^{(n)}_\mu (x;f)}$ denote the ``staple''
associated with the link $U^{(n)}_\mu(x;f)$ in terms of the entire
set of $U^{(n)}_\nu(y;f)$ matrices. One step in the recursion
takes one from a set $U^{(n)}_\mu (x;f)$ to a set $U^{(n+1)}_\mu
(x;f)$:
\begin{eqnarray}
X^{(n+1)}_\mu (x;f)= (1-|f|) U^{(n)}_\mu (x;f)+\frac{f}{6}
\Sigma_{U^{(n)}_\mu (x;f)}\nonumber\\ U^{(n+1)}_\mu (x; f
)=X^{(n+1)}_\mu (x;f) \frac{1}{\sqrt{[X^{(n+1)}_\mu (x;f)]^\dagger
X^{(n+1)}_\mu (x;f)}}
\end{eqnarray}
In the simulation, one never encounters a situation where the
unitary projection in the above equations stalls because $X^{(n)}$
is singular. In other words, smearing is well defined with
probability one.
For a twisted simulation $f$ is taken as the
negative of its value in the untwisted case; we shall quote the
conventional, positive value of $f$.

$U^{(n)}_\mu (x;f)$ transforms under gauge transformations the
same way as $U_\mu (x)$ does. Our smeared Wilson loop operators,
$\hat W [L_1,l_2 ; f ; n]$ are defined as ordered products round
the $L_1 \times L_2$ rectangle restricted to a plane. $L_\alpha$
are integers and give the size of the loop in units of the lattice
spacing. When the traversed link starts at site
$x=(x_1,x_2,x_3,x_4)$ $ x_\mu\in Z$ and connects to the
neighboring site in the positive direction $\mu$, $x+\mu$, the
link matrix is $U^{(n)}(x;f)$, while when this oriented link is
traversed in the opposite direction, the link matrix is
$U^{\dagger (n)}(x;f)$. $\hat W$ depends on the place the loop was
opened, but its eigenvalues do not. The set of eigenvalues is
gauge invariant under the fundamental gauge transformation
operating on $U_\mu (x)$.

We established by numerical means, to be described in greater
detail later, that as $L_1 , L_2 , f, n$ are varied, at specific
lattice couplings, the spectrum of $\hat W [L_1 , L_2 ,; f ; n] $
opens a gap for very small and/or very smeared loops. This gap
closes for very large and/or very lightly smeared loops.


\subsection{Lattice large $N$ transitions survive in the continuum limit.}

Even restricting to just square loops, our lattice operators
depend on many parameters: there is an overall scale, the side of
the square, $L$; there is a gauge coupling $b$ which fixes the
lattice spacing in physical units, and therefore also controls the
scale; there are two parameters, $f$ and $n$ that control the APE
smearing and are responsible of the taming of the perimeter and
corner divergences in the traces of the operator in all
irreducible $SU(N)$ representations simultaneously. Our
construction implicitly adjusts the infinite number of
undetermined finite parts associated with the renormalization of
the trace in each representation so that the entire collection can
be viewed as coming from one (fluctuating) $SU(N)$ matrix. This
ensures that backtracking segments of the loop make no
contribution to the operator.

Obviously, there is a substantial amount of arbitrariness in the
choices we made when defining the regularized Wilson loop
operators. If the construction is kept more general, this
arbitrariness would reflect itself in an invariance under a large
set of transformations in the continuum, perhaps under the form of
a generalized renormalzation group equation.

Our first objective is to find out how to adjust the parameter
dependence in $\hat W$ such that $N=\infty$ transition points
identified on the lattice, where the gap just opens, survive in
the continuum limit in which the lattice coupling $b$ is taken to
infinity together with $L_{1,2}$ in such a way that the physical
lengths $l_\alpha = L_\alpha a(b)$ are kept fixed. $a(b)$ is the
lattice spacing at coupling $b$, and we express it in terms of the
physical finite temperature deconfining transition $t_c$, by
\be
a(b)=\frac{T_c(b)}{t_c}
\ee
where $\frac{1}{T_c(b)}$ is the lattice
size of the compact direction on an $S^1\times R^3$ lattice
corresponding to the finite temperature transition critical
coupling $b$.

We set the number of smearing steps $n$ to be proportional to the
perimeter square (we restricted the loop sizes to even
$L_1+L_2$), $n=\frac{(L_1+L_2)^2}{4}$. The physical
sizes of the loop are $l_\alpha$ and $L_\alpha = \frac{l_\alpha}{a(b)}$.
We have
set $n=\frac{(L_1 (b) + L_2 (b))^2}{4}$ because in physical terms
the product $f n$ is a length squared. This can be seen from the
simple perturbative argument to be presented in
the next subsection, or, more
intuitively, by observing that smearing is a random walk that
fattens the loop and the thickness grows as the square root of the
number of smearing steps. Our choice for $n$ makes $f$ a
dimensionless parameter in the physical sense; on the lattice $f$
is actually bounded to an interval of order one.

We analyzed four sequences of $\hat W [L_1 , L_2; f;
n=\frac{(L_1 + L_2 )^2}{4}]$ keeping the physical size fixed at
four values, one for each sequence. We set $L_1=L_2=L$.
By various means, to be described later, we varied $f$ and found the
critical value of $f$ at which each member in each sequence
just opens a gap in its spectrum at eigenvalue $-1$.
(So long as CP invariance is maintained, for any loop,
the eigenvalue probability is
extremal at $\pm 1$.) The sets of critical values for
each sequence were then extrapolated
to the $b\to\infty$ limit. In each sequence
the physical size of the loop is kept fixed by definition. This gave four
continuum critical values of $f$ for the four different
physical loop sizes we selected. Throughout, the values of $f$ used were within
reasonable ranges, well within the maximally allowed lattice range
of $0<f<0.75$ \cite{pertsmear11}.
That it is actually possible to arrange that only such reasonable values of $f$
are needed in the vicinity of the large $N$ transitions
is our most basic result and a nontrivial feature of the dynamics
of planar pure Yang Mills theory. The collection of the continuum
limits of the critical values of $f$ at the
transition points associated with each sequence identified by its
physical scale defines a critical line, $f_c(l)$.
For any finite $b$ there were finite lattice spacing effects, but they
were sufficiently small to keep all the lattice approximations to
$f$ in a good range.

Thus, we conclude that we have found a continuum critical line
$f_c(l)$ in planar QCD where our continuum Wilson loop operators
have a well defined spectrum and that spectrum is critical in the
sense that its eigenvalue density, $\hat \rho (\theta)$ is nonzero
for any $\theta\ne\pi$ but $\hat\rho (\pi)=0$. Here, we denote the
eigenvalues of the Wilson loop operator as $e^{\imath\theta}$,
where $\theta$ is an angle on a circle.

It is not necessary to keep $n$ quadratic in $L_\alpha$ for
$\Lambda l_\alpha >>1$; the loops that have a critical $f$ are of
sizes $l$ for which $\Lambda l $ is of order unity. For much
larger loops, $f n$ should be allowed to become proportional to a
physical scale of order $\frac{1}{\Lambda^2}$ and not increase
further as the physical loop size goes to infinity. All that is
important is that for loop sizes of order $\frac{1}{\Lambda}$, $n$
be adjusted as above, and that for small loops we let $f n$ become
small too. For small loops we wish to enter a perturbative regime,
and therefore we should not eliminate the ultraviolet fluctuations
at short distances, where they dominate.

Once the existence of the $f_c (l)$ line has been established in
the continuum limit, with $\frac{df_c}{dl} > 0$, we can be
reasonably sure that dilating a $\hat W [L,L;f;n]$ so that at
lattice coupling $b$ it passes through a point $L=\frac{l}{a(b)},
f=f_c (l)+{\cal O}(a^2(b)), n=L^2$ the spectrum of this $\hat W$
will go through the transition. We still need to make sure that
once the gap is closed, the density at $-1$ starts building up
continuously from zero; this would mean that the transition is
continuous and likely to possess a universal character. If the
transition were discontinuous, like, for example, the finite
temperature transition in Polyakov loop operators, it is unlikely
that a scaling regime exists on either side of the transition and
a single scaling regime connecting the two sides (of the type we
are hoping for) is ruled out. This is obviously the case in the
finite temperature case because the transition there is not
induced by taking $N\to\infty$, but exists for all $N$, and starts
being first order at $N=3$.


We addressed this point by carrying out additional simulations
along another line in the physical $(f,l)$ plane, which is not a
constant $l$ but, rather, has $l$ varying by a substantial amount.
We picked a convenient interval of lengths $l$ and a line in the
plane $(f,l)$ given by $f=f_0-f_1 l^2$ with constant $f_1 >0$
chosen so that $f$ varies between $f_1$ and $f_2$ with $f_1\approx
0.1$ and $f_2\approx 0.2$. This line intersects $f_c (l)$ at an
angle not too close to zero, thus enhancing the variability in the
spectrum as one takes $\hat W$ along the line between the two
extremal lengths. There is no particular physical meaning attached
to this second line, except that along it the loop is dilated and
in addition, the amount of smearing decreased; both effects help
closing the spectral gap that is present when the loop is small.

To go along the continuum line we fixed the lattice loop-size $L$
and varied $b$. We kept $n=L^2$ also fixed, but
varied $f$ according to the formula $f=f(b)=f_0-f_1 (La(b))^2$.
Following $\hat W$
along this line we established that indeed the spectrum of $\hat
W$ opens a gap as the band of lattice curves
converging to $f_c(l)$ is crossed.

We avoided working close to the transition and estimated
independently on each side where the corresponding phase would
end, using various extrapolation methods: On the side where there
is a gap, we studied the decrease of the gap to estimate when it
will close. On the gap-less side, we estimated when $\hat \rho
(\pi)$ would vanish.

We repeated the procedure for three increasing values of $L$.
In each case one covers similar (overlapping, but not identical)
ranges of physical size $l$. In the common portions of the range,
the larger the lattice value of $L$ is, the smaller the lattice
effects ought to be. We found general consistency with
a scaling violation that goes as the lattice spacing squared.

The separate extrapolations in each phase produced independent
estimates for end points of each phase which came out
reasonably close to each other. This is consistent with a continuous large
$N$ transition at which the eigenvalue distribution of $\hat W$
opens a gap. The determinations of the critical points from each
side were ordered in a way that was inconsistent with a
discontinuous transition in which $\hat \rho(\pi)$ decreases to a
non-zero value and then drops to zero, jumping into a phase with a finite
gap around $\pi$. Thus, we conclude that the transition is
continuous and one can hope for some universal
description in the vicinity of the transition point. Actually, it
is hard to imagine what kind of dynamics could have
produced a first order large
$N$ transition in the spectrum.

\subsection{Smearing removes ultra-violet divergences.}

The effect of smearing is easy to understand in perturbation
theory where one supposes that each individual step in the
smearing iteration can be linearized. Writing $U^{(n)}_\mu
(x;f)=\exp(iA^{(n)}_\mu (x;f))$, and expanding in $A_\mu$ one
finds~\cite{pertsmear11}, in lattice Fourier space:
\be
A^{(n+1)}_\mu (q;f)= \sum_\nu h_{\mu\nu} (q) A^{(n)}_\nu (q;f) \label{eq:smearA}\ee
with
\be
h_{\mu\nu} (q)= f(q)(\delta_{\mu\nu} - \frac {{\tilde q}_\mu
{\tilde q}_\nu}{\tilde q^2})+\frac {{\tilde q}_\mu {\tilde
q}_\nu}{\tilde q^2} \label{eq:hkern}\ee where $\tilde q_\mu = 2\sin
(\frac{q_\mu}{2})$ and
\be
f(q)=1-\frac{f}{6}\tilde q^2
\label{eq:fq}
\ee
The iteration is solved by
replacing $f(q)$ by $f^n (q)$, where, for small enough $f$,
\be
f^n(q) \sim e^{-\frac{f n}{6} \tilde q^2}
\label{eq:fnsol}
\ee

The condition $0<f(q)<1$ for all momenta, requires $0<f<0.75$. To
get some feel for the numbers, a $4\times 4$ loop, at a value of
$b$ that makes it critical, will have $n=16$ and $f\sim 0.15$,
which makes $\frac{f n}{6}\sim 0.4$, which is an amount of
fattening over a distance as large as the loop itself. This is why
much larger loops should not be smeared with an $f n$ factor that
keeps on growing as the perimeter squared; rather, for a square
loop of side $L$, for example, the following choice would be
appropriate:
\be
f=\frac{\tilde f}{1+M^2 L^2}
\ee
Here, $M$ is a physical mass $m$
expressed in lattice units, $M=m a(b)$. $m$ is a typical hadronic
scale chosen so that at the large $N$ transition we found, $ML$ is
less then $0.01$, say.

To reassure ourselves that such an extension of the cutoff
dependence of $f$ indeed would connect smoothly to large loops
which show confinement with the known value of the string tension,
we roughly estimated the string tension from our smeared loops,
at the large size end of the ranges we studied. We found that the
string tension comes out close to that of bare lattice Wilson
loops, and this shows that our smeared Wilson loop operators are
acceptable pure $SU(N)$ gauge theory string candidates.


In summary, the effect of smearing on the propagator of
$A^{(n)}_\mu (x;f)$ is to eliminate the ultraviolet singularity
that gives rise to perimeter and corner divergences in traces of
bare Wilson loop operators. This works when the product $fn$
varies as a physical length squared when the continuum limit is
approached. For small loops we choose this physical length to be
proportional to the loop itself. For large loops the physical
length stabilizes at some hadronic scale. In this way, our
operators become sensitive to the short distance dynamics when the
scale is very small, but, when confinement becomes the dominating
effect, for very large loops, the string tension of the smeared
Wilson loops is the true, universal, string tension of the theory,
unaffected by the smearing. The leading term in the effective
string theory, describing the shape dependence of the Wilson loops
when the loop is very large, is expected to be universal and the
string tension enters only in a trivial manner, as a dimensional
unit.

We stress that we are not executing a renormalization group
transformation; our Wilson loop operators are defined for
infinitely thin loops in space-time and for all loop sizes;
rather, we are applying a regularization whose objective is to
make the Wilson loop operators finite in the continuum. If we
dealt with each irreducible representation separately, the
elimination of perimeter and corner divergences would leave behind
an infinite number of arbitrary finite parts. This freedom has
been implicitly exploited by imposing the requirement of unitarity
of the Wilson loop operators, which ensures what Polyakov has
termed the ``zigzag'' symmetry \cite{poly3}. Preserving the
unitary character of the Wilson loop operator in turn makes the
concept of an eigenvalue distribution well defined in the
continuum and sets the stage for possible large $N$ phase
transitions. The zigzag symmetry also influences the form of the
lattice loop equations at infinite $N$, which are obeyed by traces
of bare Wilson loops in the fundamental representation; the
effective string theory for large loops needs not obey the zigzag
symmetry, but this symmetry could provide a useful constraint on
its couplings.
